\documentclass[11pt,a4paper]{article}

%  XeLaTeX
\usepackage{fontspec}

\usepackage{algorithm}
\usepackage{algorithmic}
\usepackage{amsmath,amssymb,amsthm}
\usepackage{bm}
\usepackage[useregional]{datetime2}
\usepackage{enumitem}
\usepackage{booktabs}
\usepackage{float}
\usepackage{graphicx}
\usepackage{tabularray}
\usepackage{framed}
\usepackage[a4paper, margin=2.5cm]{geometry}
\usepackage{eso-pic}
\usepackage{mathtools}
\usepackage{tikz}
\usetikzlibrary{calc}
\usepackage{xcolor}

\usepackage[backend=biber, style=apa, sorting=nyt]{biblatex}
\addbibresource{D:/github/ENEL445/report/references.bib}
\graphicspath{{D:/github/ENEL445/figs/}}

\usepackage[colorlinks=true, linkcolor=blue, citecolor=blue, urlcolor=blue]{hyperref}
\hypersetup{
  pdfauthor   = {David Ewing},
  pdfsubject  = {\jobname.tex},
  pdfkeywords = {source: \jobname.tex}
}

% Running margin label
\newcommand{\mymarginlabel}{\textcopyright\ David Ewing dew59@uclive.ac.nz | \DTMnow\ | \jobname.tex}
\AddToShipoutPictureBG{%
  \begin{tikzpicture}[remember picture,overlay]
    \node[rotate=90, anchor=north]
      at ($(current page.north west)+(1.4cm,-13cm)$)
      {\scriptsize\ttfamily \mymarginlabel};
  \end{tikzpicture}%
}

% Custom commands
\newcommand{\E}{\mathbb{E}}
\newcommand{\R}{\mathbb{R}}
\newcommand{\N}{\mathcal{N}}
\newcommand{\ELBO}{\text{ELBO}}
\newcommand{\tr}{\text{tr}}
\newcommand{\sigmoid}{\sigma}
\setlength{\parindent}{0pt}

\title{
Case Study 3 --- Bayesian Logistic Regression:\\[0.4em]
Gradient-Based Optimisation and Posterior Variance Collapse
}
\author{
  David Ewing (82171165)
}
\date{9 May 2026}

\begin{document}

\maketitle

\begin{abstract}
\setlength{\parindent}{0pt}
\setlength{\parskip}{0.7em}

This report applies the gradient-based optimisation methods from ENEL445 to variational
inference in Bayesian logistic regression.
Unlike the Gaussian regression setting, logistic regression has no conjugate variational
family, and the Evidence Lower Bound (ELBO) does not admit a simple closed-form update.
The Jaakkola--Jordan \parencite{Jaakkola2000} variational lower bound provides a tractable
ELBO for logistic regression via the local variational bound $\log\sigma(\psi) \geq
\tfrac{\psi}{2} - \log 2\cosh(\xi/2)$, which is tight when $\xi = \sqrt{\E[\psi^2]}$.

Four optimisation methods are applied to maximise the ELBO:
Coordinate Ascent Variational Inference (CAVI), gradient ascent with Armijo backtracking,
Newton's method with regularised finite-difference Hessian, and BFGS with Wolfe conditions.
The variational family is $q(\bm{\beta}) = \N(\bm{\mu},\bm{\Sigma})$ with a
Cholesky reparameterisation to preserve positive definiteness throughout optimisation.

The reference posterior is obtained from a P\'{o}lya--Gamma Gibbs sampler
\parencite{Polson2013} using a truncated sum-of-gammas approximation with $T=50$ terms.
Diagnostics include effective sample size (ESS), autocorrelation function (ACF) plots,
and computational cost per effective sample.

The model is $y_i \sim \mathrm{Bernoulli}(\sigma(\beta_0 + \beta_1 x_i))$
with $\bm{\beta}_{\mathrm{true}} = (-1.0,\; 2.0)$ and $n=200$ observations.
All four VI methods recover posterior means within $0.02$ of the Gibbs reference.
CAVI converges in fewer than 50 iterations and is several orders of magnitude faster
than the P\'{o}lya--Gamma sampler, which requires $200$ Gamma draws per iteration.
Note that timing comparisons across the linear, quadratic, and logistic notebooks
in this project are not controlled experiments: the likelihoods, parameter counts,
and Gibbs structures differ across cases.
The fair within-notebook comparison is each VI method versus its own Gibbs reference.
\end{abstract}

\newpage

%─────────────────────────────────────────────────────────────────────────────
\section*{Introduction}
%─────────────────────────────────────────────────────────────────────────────

Bayesian inference provides principled uncertainty quantification for binary classification
through the posterior distribution
\begin{equation*}
  p(\bm{\beta} \mid \bm{y}) \propto p(\bm{y} \mid \bm{\beta})\,p(\bm{\beta}),
\end{equation*}
where $p(\bm{y}\mid\bm{\beta}) = \prod_{i=1}^{n} \sigma(x_i^T\bm{\beta})^{y_i}
(1-\sigma(x_i^T\bm{\beta}))^{1-y_i}$ is the Bernoulli log-likelihood under the logistic
link $\sigma(\psi)=(1+e^{-\psi})^{-1}$.
The posterior is not available in closed form because the logistic likelihood is not
conjugate to the Gaussian prior $p(\bm{\beta})=\N(\bm{0}, \lambda^{-1}\bm{I})$.

Variational Inference (VI) converts posterior computation into optimisation by finding the
best approximation $q(\bm{\beta})$ from a tractable family, maximising the Evidence Lower
Bound (ELBO):
\begin{equation*}
  \mathcal{L}(\bm{\phi}) = \E_{q}[\log p(\bm{y}\mid\bm{\beta})] - \mathrm{KL}(q \| p)
  \leq \log p(\bm{y}).
\end{equation*}
The logistic likelihood $\log\sigma(\psi)$ is not quadratic in $\bm{\beta}$, so the
expectation $\E_q[\log\sigma(x_i^T\bm{\beta})]$ cannot be computed analytically under
a Gaussian $q$.
The Jaakkola--Jordan bound replaces this term with a tight quadratic lower bound,
restoring conjugacy and making CAVI and gradient-based methods tractable.

This report is the third numerical study in this project (following linear and quadratic
Gaussian regression), and addresses a qualitatively different challenge: the logistic
likelihood introduces a non-conjugate computational bottleneck that gradient-based methods
must navigate.

%─────────────────────────────────────────────────────────────────────────────
\section{Model Formulation}
%─────────────────────────────────────────────────────────────────────────────

\subsection{Bayesian Logistic Regression}

The model is
\begin{align}
  y_i \mid \bm{\beta} &\sim \mathrm{Bernoulli}(\sigma(x_i^T\bm{\beta})), \quad i=1,\ldots,n,\\
  \bm{\beta} &\sim \N(\bm{0}, \lambda^{-1}\bm{I}_p), \quad \lambda^{-1}=10,
\end{align}
where $x_i = (1, x_i)^T \in \R^2$ is the feature vector (intercept + covariate),
and $\sigma(\psi) = (1+e^{-\psi})^{-1}$ is the logistic sigmoid.

The test case uses $n=200$ observations with $x_i \sim \mathrm{Uniform}(-2,2)$ and
true parameters $\bm{\beta}_{\mathrm{true}} = (-1.0,\; 2.0)$.

\subsection{Mean-Field Variational Approximation}

The variational family is
\begin{align}
  q(\bm{\beta}) &= \N(\bm{\mu}, \bm{\Sigma}),
\end{align}
with variational parameters $\bm{\phi} = (\bm{\mu}, \bm{\Sigma})$.
The covariance $\bm{\Sigma} = \bm{L}\bm{L}^T$ is represented by its Cholesky factor
$\bm{L}$ to ensure positive definiteness throughout optimisation.

The unconstrained parameter vector is
\begin{equation}
  \bm{\theta} = (\mu_0,\; \mu_1,\; \tilde{\ell}_{11},\; \ell_{21},\; \tilde{\ell}_{22})
  \in \R^5,
\end{equation}
where $L_{11} = e^{\tilde{\ell}_{11}}$ and $L_{22} = e^{\tilde{\ell}_{22}}$ enforce
$L_{ii} > 0$.

\subsection{Jaakkola--Jordan ELBO}

The log-likelihood term $\E_q[\log\sigma(\psi_i)]$ where $\psi_i = x_i^T\bm{\beta}$
is intractable under a Gaussian $q$.
The Jaakkola--Jordan bound uses a local variational parameter $\xi_i > 0$ to introduce
a quadratic lower bound:
\begin{equation}
  \log\sigma(\psi_i) \geq \frac{\psi_i}{2} - \log 2\cosh\!\left(\frac{\xi_i}{2}\right)
  - \frac{\xi_i}{2}\bigl(\E_q[\psi_i^2]/\xi_i^2 - 1\bigr) \cdot \frac{\tanh(\xi_i/2)}{4}.
\end{equation}
At the optimal $\xi_i^* = \sqrt{\E_q[\psi_i^2]}$, where
$\E_q[\psi_i^2] = (x_i^T\bm{\mu})^2 + x_i^T\bm{\Sigma} x_i$,
the bound is tight and the ELBO becomes
\begin{equation}\label{eq:elbo}
  \mathcal{L}(\bm{\theta}) = \sum_{i=1}^{n}
  \Bigl[(y_i - \tfrac{1}{2})\,x_i^T\bm{\mu}
        - \log 2 - \log\cosh\!\tfrac{\xi_i^*}{2}\Bigr]
  - \mathrm{KL}(q\|p),
\end{equation}
where the KL divergence for Gaussian families is
\begin{equation}
  \mathrm{KL}(q\|p) = \tfrac{1}{2}\Bigl[
    \log\frac{|\bm{\Sigma}_0|}{|\bm{\Sigma}|}
    - p + \tr(\bm{\Sigma}_0^{-1}\bm{\Sigma})
    + (\bm{\mu}-\bm{\mu}_0)^T\bm{\Sigma}_0^{-1}(\bm{\mu}-\bm{\mu}_0)
  \Bigr],
\end{equation}
with $\bm{\mu}_0=\bm{0}$ and $\bm{\Sigma}_0 = 10\bm{I}$.

\subsection{Key Notation}

\begin{center}
\begin{tabular}{@{}ll@{}}
\toprule
Symbol & Meaning \\
\midrule
\multicolumn{2}{@{}l}{\textbf{Bayesian model}} \\
$\bm{X}$ & Design matrix ($n \times p$), rows $x_i^T$ \\
$\bm{y}$ & Binary response vector ($y_i \in \{0,1\}$) \\
$\bm{\beta}$ & Regression coefficient vector \\
$\sigma(\cdot)$ & Logistic sigmoid $\sigma(\psi)=(1+e^{-\psi})^{-1}$ \\
$n, p$ & Observations, predictors \\
\midrule
\multicolumn{2}{@{}l}{\textbf{Variational approximation}} \\
$q(\bm{\beta})$ & Variational Gaussian approximation \\
$\bm{\mu}, \bm{\Sigma}$ & Mean and covariance of $q(\bm{\beta})$ \\
$\bm{L}$ & Cholesky factor, $\bm{\Sigma}=\bm{L}\bm{L}^T$ \\
$\mathcal{L}(\bm{\theta})$ & Jaakkola--Jordan ELBO \\
$\xi_i$ & Local variational parameter (tightness parameter) \\
\midrule
\multicolumn{2}{@{}l}{\textbf{Optimisation}} \\
$\bm{\theta}$ & Unconstrained vector $(\bm{\mu}, \tilde{\ell}_{11}, \ell_{21}, \tilde{\ell}_{22})$ \\
$\bm{d}^{(t)}$ & Search direction at iteration $t$ \\
$\alpha_t$ & Step size at iteration $t$ \\
\midrule
\multicolumn{2}{@{}l}{\textbf{P\'{o}lya--Gamma sampler}} \\
$\omega_i$ & P\'{o}lya--Gamma latent variable, $\omega_i \mid \bm{\beta} \sim \mathrm{PG}(1, x_i^T\bm{\beta})$ \\
$T$ & Truncation order for PG series approximation \\
$\kappa_i$ & Pseudo-response $\kappa_i = y_i - 1/2$ \\
$\mathrm{ESS}$ & Effective sample size \\
\midrule
\multicolumn{2}{@{}l}{\textbf{Lecture correspondence} (Prof.\ R.\,S.\ Smith, ENEL445 2026)} \\
$J(\theta)$ & Cost function $=\mathcal{L}(\bm{\theta})$ (ELBO) in this report \\
$\Phi$ & Regressor matrix $=\bm{X}$ in this report \\
\bottomrule
\end{tabular}
\end{center}

%─────────────────────────────────────────────────────────────────────────────
\section{P\'{o}lya--Gamma Gibbs Sampler}
%─────────────────────────────────────────────────────────────────────────────

The P\'{o}lya--Gamma (PG) Gibbs sampler of \textcite{Polson2013} provides a reference
posterior for assessing variational approximation quality.
The key identity is
\begin{equation}
  \frac{(e^\psi)^a}{(1+e^\psi)^b}
  = 2^{-b} e^{\kappa\psi}
    \int_0^\infty e^{-\omega\psi^2/2}\, p(\omega)\,\mathrm{d}\omega,
  \quad \omega \sim \mathrm{PG}(b, 0),
\end{equation}
which introduces a P\'{o}lya--Gamma latent variable $\omega_i$ that renders the
logistic likelihood conditionally Gaussian.

\subsection{Gibbs Updates}

The joint model with PG augmentation is
\begin{align}
  \bm{\beta} \mid \bm{\omega}, \bm{y} &\sim \N(m_\omega, V_\omega),\\
  \omega_i \mid \bm{\beta} &\sim \mathrm{PG}(1,\, x_i^T\bm{\beta}),
\end{align}
where
\begin{align}
  V_\omega &= \bigl(\bm{X}^T \mathrm{diag}(\bm{\omega})\bm{X} + \bm{\Sigma}_0^{-1}\bigr)^{-1},\\
  m_\omega &= V_\omega\bigl(\bm{X}^T\bm{\kappa} + \bm{\Sigma}_0^{-1}\bm{\mu}_0\bigr),
  \quad \kappa_i = y_i - \tfrac{1}{2}.
\end{align}
Both full conditionals are available analytically, so each Gibbs iteration is exact
(up to the PG truncation).

\subsection{PG Sampling Approximation}

The PG distribution $\mathrm{PG}(1, c)$ is sampled using the truncated series
\begin{equation}
  \omega \stackrel{d}{=} \frac{1}{2\pi^2}
  \sum_{k=1}^{T} \frac{g_k}{(k-\tfrac{1}{2})^2 + c^2/(4\pi^2)},
  \quad g_k \sim \mathrm{Gamma}(1,1),
\end{equation}
with $T=50$ terms.
This is sufficient for $|\psi_i| \lesssim 4$ with negligible truncation error.
The sampler runs for $N_{\mathrm{warmup}} = 500$ burn-in iterations and
$N = 2000$ retained samples.

\textbf{Note on timing:}
Each PG Gibbs iteration requires $n \times T = 200 \times 50 = 10\,000$ Gamma variates.
This is substantially more expensive per sample than the conjugate Gaussian Gibbs used in
the linear and quadratic cases, so cross-model timing comparisons are qualitative only.

%─────────────────────────────────────────────────────────────────────────────
\section{Optimisation Problem Formulation}
%─────────────────────────────────────────────────────────────────────────────

\subsection*{Decision Variables}

The optimisation is carried out over the unconstrained vector
\begin{equation}
  \bm{\theta} = (\mu_0,\; \mu_1,\; \tilde{\ell}_{11},\; \ell_{21},\; \tilde{\ell}_{22})
  \in \R^5,
\end{equation}
so that $\bm{\mu} = (\mu_0, \mu_1)^T$ and
$\bm{\Sigma} = \bm{L}\bm{L}^T$ with $\bm{L} = \bigl[\begin{smallmatrix}
e^{\tilde{\ell}_{11}} & 0 \\ \ell_{21} & e^{\tilde{\ell}_{22}} \end{smallmatrix}\bigr]$.
The constraint $\bm{\Sigma} \succ 0$ is satisfied automatically by this parameterisation.

\subsection*{Objective Function}

Maximise the Jaakkola--Jordan ELBO (Equation~\ref{eq:elbo}) at the optimal
$\xi_i^* = \sqrt{(x_i^T\bm{\mu})^2 + x_i^T\bm{\Sigma} x_i}$.
The function $\mathcal{L}: \R^5 \to \R$ is smooth, bounded above by $\log p(\bm{y})$,
and has a unique stationary point in this model (the CAVI fixed point).

\subsection*{Constraints}

No explicit constraints remain after reparameterisation.
The Cholesky parameterisation maps $\R^5$ bijectively to the space of
valid Gaussian variational parameters.

\subsection*{Optimisation Components}

Each method is described in terms of the four components from the ENEL445 course:
entry conditions, search direction, step-size rule, and exit conditions.

\begin{description}
\item[CAVI.]
  \textit{Entry:} always applicable.
  \textit{Direction:} closed-form coordinate update $\bm{\Sigma}^{-1} \leftarrow
  \bm{\Sigma}_0^{-1} + 2\sum_i\lambda(\xi_i^*)x_ix_i^T$,\;
  $\bm{\mu} \leftarrow \bm{\Sigma}(\bm{\Sigma}_0^{-1}\bm{\mu}_0 + \bm{X}^T\bm{\kappa})$,
  where $\lambda(\xi) = \tanh(\xi/2)/(4\xi)$.
  \textit{Step size:} unit (exact closed-form step).
  \textit{Exit:} $|\Delta\mathcal{L}| < \varepsilon$.

\item[Gradient ascent.]
  \textit{Entry:} $\nabla\mathcal{L}(\bm{\theta}^{(t)})^T\bm{d}^{(t)} > 0$ (direction is
  the gradient itself; always satisfied).
  \textit{Direction:} $\bm{d}^{(t)} = \nabla\mathcal{L}(\bm{\theta}^{(t)})$.
  \textit{Step size:} Armijo backtracking, starting from $\alpha_0 = 0.5$, halving until
  $\mathcal{L}(\bm{\theta}+\alpha\bm{d}) \geq \mathcal{L}(\bm{\theta}) + c_1\alpha\|\bm{g}\|^2$,
  $c_1 = 10^{-4}$.
  \textit{Exit:} $\|\bm{g}\|_2 < 10^{-8}$ or $t \geq 2000$.

\item[Newton's method.]
  \textit{Entry:} ascent condition $\bm{g}^T\bm{d} > 0$ enforced by Hessian regularisation
  $\bm{H}_{\mathrm{reg}} = \bm{H} - \lambda\bm{I}$ with $\lambda$ increased until
  $-\bm{H}_{\mathrm{reg}}^{-1}\bm{g}$ is an ascent direction.
  \textit{Direction:} $\bm{d}^{(t)} = -\bm{H}_{\mathrm{reg}}^{-1}\bm{g}^{(t)}$ via finite-difference
  Hessian.
  \textit{Step size:} Armijo backtracking.
  \textit{Exit:} $\|\Delta\bm{\theta}\| < 10^{-8}$ or $t \geq 100$.

\item[BFGS.]
  \textit{Entry:} quasi-Newton direction satisfies ascent condition when $\bm{y}_t^T\bm{s}_t>0$.
  Wolfe conditions ensure this curvature condition is maintained.
  \textit{Direction:} $\bm{d}^{(t)} = -\bm{B}_t \bm{g}^{(t)}$ where $\bm{B}_t$ is the
  inverse-Hessian approximation (positive definite, updated via BFGS rank-2 formula).
  \textit{Step size:} Wolfe conditions (scipy implementation).
  \textit{Exit:} $\|\bm{g}\| < 10^{-8}$ or $t \geq 2000$.
\end{description}

%─────────────────────────────────────────────────────────────────────────────
\section{Numerical Study}
%─────────────────────────────────────────────────────────────────────────────

\subsection{Test Case Design}

\begin{itemize}
  \item Model: $y_i \sim \mathrm{Bernoulli}(\sigma(\beta_0 + \beta_1 x_i))$
  \item True parameters: $\bm{\beta}_{\mathrm{true}} = (-1.0,\; 2.0)$
  \item Covariates: $x_i \sim \mathrm{Uniform}(-2, 2)$, $n=200$
  \item Prior: $\bm{\beta} \sim \N(\bm{0}, 10\bm{I})$
  \item Seed: \texttt{82171165} (student ID)
\end{itemize}

\subsection{Gradient Verification}

Before running any optimiser, the analytical gradient of the ELBO with respect to
$\bm{\theta}$ was verified against central finite differences at a non-trivial
starting point $\bm{\theta}_0 = (\bm{0},\, \log 0.1, 0, \log 0.1)$.
All components agreed to relative error below $10^{-6}$.

\begin{figure}[H]
  \centering
  \includegraphics[width=0.80\linewidth]{logistic_elbo_gradient_check.png}
  \caption{Analytical versus finite-difference gradient at $\bm{\theta}_0$.
    Bars coincide visually for all five components, confirming correct gradient
    implementation.}
  \label{fig:gradcheck}
\end{figure}

\subsection{ELBO Landscape}

Figure~\ref{fig:landscape} shows a one-dimensional slice of the ELBO surface
along $\mu_1$ with all other parameters held at $\bm{\theta}_0$.
The surface is smooth and unimodal, with the maximum near $\beta_{1,\mathrm{true}}=2.0$.

\begin{figure}[H]
  \centering
  \includegraphics[width=0.60\linewidth]{logistic_elbo_landscape.png}
  \caption{ELBO slice along $\mu_1$ at $\bm{\theta}_0$. The dashed red line marks the
    true value $\beta_1 = 2.0$.}
  \label{fig:landscape}
\end{figure}

%─────────────────────────────────────────────────────────────────────────────
\section{Results}
%─────────────────────────────────────────────────────────────────────────────

\subsection{Generated Data}

Figure~\ref{fig:data} shows the generated binary dataset and the true sigmoid curve.
Of the $n=200$ observations, approximately 105 have $y_i=1$.

\begin{figure}[H]
  \centering
  \includegraphics[width=0.75\linewidth]{logistic_data_scatter.png}
  \caption{Generated logistic data ($n=200$). The solid line is the true sigmoid
    $\sigma(\beta_0 + \beta_1 x)$.}
  \label{fig:data}
\end{figure}

\subsection{ELBO Convergence}

Figure~\ref{fig:elbocv} shows the ELBO trajectory for CAVI, gradient ascent, and
Newton's method (BFGS is not shown as its trajectory is internal to scipy).
CAVI converges monotonically in fewer than 50 iterations.
Gradient ascent converges more slowly but reaches the same final value as CAVI when
given sufficient iterations.
Newton's method converges rapidly in terms of iteration count, but each iteration is
more expensive due to the finite-difference Hessian computation.

\begin{figure}[H]
  \centering
  \includegraphics[width=\linewidth]{logistic_elbo_convergence.png}
  \caption{ELBO convergence (left) and convergence rate on a log scale (right) for CAVI,
    gradient ascent, and Newton's method.}
  \label{fig:elbocv}
\end{figure}

\subsection{Step-Size Behaviour}

Figure~\ref{fig:stepsize} shows the Armijo step sizes for gradient ascent and
Newton's method.
Gradient ascent uses a consistently small step size throughout, reflecting the
slowly varying curvature of the Jaakkola--Jordan ELBO.
Newton's method uses unit step size for most iterations (the Armijo condition is
satisfied on the first trial), consistent with near-quadratic convergence near the optimum.

\begin{figure}[H]
  \centering
  \includegraphics[width=\linewidth]{logistic_gradient_stepsize.png}
  \caption{Armijo step size per iteration for gradient ascent (left) and Newton's method
    (right). Newton achieves near-unit steps near convergence.}
  \label{fig:stepsize}
\end{figure}

\subsection{Newton Hessian Condition Number}

Figure~\ref{fig:cond} shows the condition number of the regularised Hessian used in
Newton's method.
Early iterations may require significant regularisation ($\lambda \gg 0$) if the
finite-difference Hessian is ill-conditioned; near convergence the Hessian becomes
better conditioned and regularisation is minimal.

\begin{figure}[H]
  \centering
  \includegraphics[width=0.65\linewidth]{logistic_newton_condition.png}
  \caption{Condition number of the regularised Hessian at each Newton iteration.}
  \label{fig:cond}
\end{figure}

\subsection{P\'{o}lya--Gamma Gibbs Sampler Diagnostics}

Figures~\ref{fig:trace} and~\ref{fig:acf} show the trace plots and ACF for the two
parameters from the PG Gibbs sampler.
The chains mix well, with ACF dropping below the 95\% confidence band within 10--20 lags.

\begin{figure}[H]
  \centering
  \includegraphics[width=\linewidth]{logistic_gibbs_trace_beta0.png}
  \caption{Trace plot (left) and marginal posterior histogram (right) for $\beta_0$
    from the PG Gibbs sampler.  Red dashed: true value; black dotted: posterior mean.}
  \label{fig:trace}
\end{figure}

\begin{figure}[H]
  \centering
  \includegraphics[width=\linewidth]{logistic_gibbs_trace_beta1.png}
  \caption{Trace plot and posterior histogram for $\beta_1$.}
  \label{fig:trace1}
\end{figure}

\begin{figure}[H]
  \centering
  \includegraphics[width=\linewidth]{logistic_gibbs_acf_beta0.png}
  \caption{Autocorrelation function for $\beta_0$ (PG Gibbs). Red dashed lines are
    the 95\% confidence band for white noise.}
  \label{fig:acf}
\end{figure}

\begin{figure}[H]
  \centering
  \includegraphics[width=\linewidth]{logistic_gibbs_acf_beta1.png}
  \caption{Autocorrelation function for $\beta_1$.}
  \label{fig:acf1}
\end{figure}

\subsection{Effective Sample Size and Cost per ESS}

Figure~\ref{fig:ess} reports the ESS for both parameters and the computational cost
per effective sample.
With $T=50$ PG truncation terms and $N=2000$ samples, ESS exceeds 1000 for both
parameters, indicating good mixing.

\begin{figure}[H]
  \centering
  \includegraphics[width=\linewidth]{logistic_gibbs_ess.png}
  \caption{Left: effective sample size for each parameter. Right: cost per effective
    sample (ms). VI methods are optimisers, not samplers; cost-per-ESS applies only
    to the Gibbs reference.}
  \label{fig:ess}
\end{figure}

\subsection{Posterior Comparison: VI versus Gibbs}

Figures~\ref{fig:vbgibbs0} and~\ref{fig:vbgibbs1} overlay the Gaussian VI approximations
against the PG Gibbs histogram for each parameter.
All four VI methods agree closely with the Gibbs reference for both posterior means and
standard deviations, indicating minimal posterior variance collapse in this setting.
This is consistent with the Jaakkola--Jordan bound being tight at the CAVI optimum.

\begin{figure}[H]
  \centering
  \includegraphics[width=0.80\linewidth]{logistic_vb_gibbs_beta0.png}
  \caption{Posterior comparison for $\beta_0$: Gibbs histogram versus VI Gaussian
    approximations (CAVI, gradient ascent, BFGS). Red dashed: true value.}
  \label{fig:vbgibbs0}
\end{figure}

\begin{figure}[H]
  \centering
  \includegraphics[width=0.80\linewidth]{logistic_vb_gibbs_beta1.png}
  \caption{Posterior comparison for $\beta_1$.}
  \label{fig:vbgibbs1}
\end{figure}

\subsection{Predictive Sigmoid}

Figure~\ref{fig:pred} shows the estimated sigmoid curve from BFGS and Gibbs posterior
mean compared to the true curve.
Both estimates are visually indistinguishable from the truth over the range of the
training data, confirming that both methods accurately recover the classification boundary.

\begin{figure}[H]
  \centering
  \includegraphics[width=0.75\linewidth]{logistic_predictive.png}
  \caption{Estimated predictive sigmoid $\sigma(x^T\hat{\bm{\beta}})$ for BFGS (VI)
    and Gibbs posterior mean versus the true curve.}
  \label{fig:pred}
\end{figure}

\subsection{Posterior Standard-Deviation Ratios}

Figure~\ref{fig:sdratios} shows the ratio $\sigma_{\mathrm{VI}} / \sigma_{\mathrm{Gibbs}}$
for each method and parameter.
Ratios below 1 indicate posterior variance collapse.
All methods achieve ratios between 0.90 and 1.10, indicating minimal collapse in
this two-parameter logistic model with $n=200$ observations.

\begin{figure}[H]
  \centering
  \includegraphics[width=0.80\linewidth]{logistic_sd_ratios.png}
  \caption{Posterior standard-deviation ratios (VI / Gibbs) for each method and parameter.
    The dashed line marks the ideal ratio of 1.0. Values below 1 indicate collapse.}
  \label{fig:sdratios}
\end{figure}

\subsection{Posterior Mean Errors}

Figure~\ref{fig:meanerr} shows the absolute error $|\mu_{\mathrm{VI}} - \mu_{\mathrm{Gibbs}}|$
for each method and parameter.
All methods recover posterior means to within 0.02, confirming reliable mean estimation
regardless of which VI method is used.

\begin{figure}[H]
  \centering
  \includegraphics[width=0.80\linewidth]{logistic_mean_errors.png}
  \caption{Absolute posterior mean errors relative to the Gibbs reference.}
  \label{fig:meanerr}
\end{figure}

\subsection{Full Posterior Comparison: All Methods}

Figures~\ref{fig:comp0} and~\ref{fig:comp1} show the full overlay of all four VI methods
plus the Gibbs reference for each parameter.

\begin{figure}[H]
  \centering
  \includegraphics[width=0.80\linewidth]{logistic_comparison_beta0.png}
  \caption{All methods --- posterior distribution for $\beta_0$. The Gibbs histogram
    (grey) is the reference; coloured curves are the VI Gaussian approximations.}
  \label{fig:comp0}
\end{figure}

\begin{figure}[H]
  \centering
  \includegraphics[width=0.80\linewidth]{logistic_comparison_beta1.png}
  \caption{All methods --- posterior distribution for $\beta_1$.}
  \label{fig:comp1}
\end{figure}

\subsection{Timing Comparison}

Figure~\ref{fig:timing} shows the computational cost of each method on a log scale,
with annotations giving the speedup factor relative to the PG Gibbs sampler.

\begin{figure}[H]
  \centering
  \includegraphics[width=0.90\linewidth]{logistic_timing_comparison.png}
  \caption{Computational cost on a log scale. Annotations give speedup relative to
    the PG Gibbs reference (dashed line). Note that cross-model timing comparisons
    (vs.\ the linear and quadratic cases) are qualitative only: the PG sampler
    requires $n \times T = 10\,000$ Gamma variates per iteration.}
  \label{fig:timing}
\end{figure}

%─────────────────────────────────────────────────────────────────────────────
\section{Discussion}
%─────────────────────────────────────────────────────────────────────────────

\subsection{Posterior Variance Collapse}

Unlike the linear and quadratic Gaussian regression cases, the logistic model with the
Jaakkola--Jordan bound does not exhibit strong posterior variance collapse.
This is because the bound is tight at the CAVI fixed point: when $\xi_i^* = \sqrt{\E_q[\psi_i^2]}$,
the bound equals the true expectation, and the ELBO objective penalises over-concentration.
The result is that CAVI and all gradient-based methods recover both posterior means and
posterior standard deviations accurately relative to the PG Gibbs reference.

\subsection{Optimisation Behaviour}

CAVI is the most efficient method due to its closed-form updates under the
Jaakkola--Jordan bound.
Gradient ascent converges more slowly but requires only gradient computations.
Newton's method converges in fewer iterations but requires a finite-difference Hessian
($25$ evaluations per iteration), making each step expensive.
BFGS provides a good compromise: quasi-Newton curvature information with only gradient
evaluations, and converges rapidly due to the smooth, near-quadratic ELBO landscape.

\subsection{Comparison Across Test Cases}

\begin{center}
\begin{tabular}{@{}llll@{}}
\toprule
Property & Linear & Quadratic & Logistic (this report) \\
\midrule
Likelihood & Gaussian & Gaussian & Bernoulli \\
Conjugacy & Yes & Yes & No (JJ bound) \\
Reference sampler & Gaussian Gibbs & Gaussian Gibbs & PG Gibbs \\
Gibbs cost & Low & Low & High ($n \times T$ Gammas) \\
Variance collapse & Mild & Moderate & Minimal \\
CAVI updates & Closed form & Closed form & Closed form (via JJ) \\
\bottomrule
\end{tabular}
\end{center}

Note that all timing figures across the three test cases are from the same machine
(Alienware 17~R5, Intel i7-8750H) but are not controlled comparisons.
Model complexity, parameter count, and Gibbs overhead differ across cases.

%─────────────────────────────────────────────────────────────────────────────
\section*{Conclusion}
%─────────────────────────────────────────────────────────────────────────────

This report has applied gradient-based optimisation methods to variational inference
in Bayesian logistic regression using the Jaakkola--Jordan ELBO.
All four VI methods (CAVI, gradient ascent, Newton, BFGS) successfully maximise the ELBO
and recover posterior distributions that closely match the PG Gibbs reference.
Posterior variance collapse is minimal in this setting, in contrast to the linear and
quadratic Gaussian cases.

The main practical finding is that CAVI is the preferred method for this model:
it converges rapidly via closed-form coordinate updates derived from the JJ bound,
requires no hyperparameter tuning, and achieves the same accuracy as more
expensive gradient-based methods.
BFGS is the recommended gradient-based alternative when closed-form updates are
not available.

The PG Gibbs sampler provides an accurate reference but is computationally expensive
due to the $n \times T$ Gamma draws required per iteration.
For production logistic regression with $n \gg 200$, the VI methods offer a substantial
computational advantage.

\printbibliography

%─────────────────────────────────────────────────────────────────────────────
\appendix
\section{Code Structure}
\label{app:code}
%─────────────────────────────────────────────────────────────────────────────

The implementation is in \texttt{python/logistic\_run.ipynb} (35 cells).

\begin{description}
  \item[Stage 1 (cells 2--4):] Data generation.
    Saves \texttt{results/data/logistic\_data.parquet}.
  \item[Stage 2 (cells 5--9):] ELBO evaluation and gradient verification.
    Functions \texttt{elbo()}, \texttt{elbo\_grad()}, \texttt{unpack()}, \texttt{pack()}.
  \item[Stage 3 (cells 10--18):] Optimisation methods.
    Functions \texttt{run\_cavi()}, \texttt{run\_gradient\_ascent()},
    \texttt{run\_newton()}, BFGS via \texttt{scipy.optimize.minimize}.
    Saves \texttt{results/data/logistic\_opt.parquet}.
  \item[Stage 4 (cells 19--24):] PG Gibbs sampler with ESS diagnostics.
    Functions \texttt{sample\_pg\_one()}, \texttt{sample\_pg\_vec()},
    \texttt{acf()}, \texttt{compute\_ess()}.
    Saves \texttt{results/data/logistic\_gibbs.parquet}.
  \item[Stage 5 (cells 25--34):] Diagnostics and comparison.
    Saves \texttt{results/tables/logistic\_diagnostics.csv} and \texttt{.parquet}.
\end{description}

All figures are saved to \texttt{figs/} with the \texttt{logistic\_} prefix to avoid
collisions with the \texttt{linear\_} and \texttt{quadratic\_} figures from the
other test cases.

\end{document}
